% =============================================================================
% CHAPTER 02 -- E-Commerce Cohort Analysis (Olist)
% =============================================================================

\chapter{E-Commerce Cohort Analysis}
\label{ch:ecommerce}

\sourcefiles{%
  \filepath{projects/02-ecommerce-cohort-analysis/} ---
  notebooks (4), SQL scripts (5), Streamlit app (6 pages)
}

This chapter documents the analytical work behind the Olist Brazilian E-Commerce
cohort analysis. The central business question is deceptively simple: with only
3\% of customers making a repeat purchase, what distinguishes the returning
minority? The project delivers statistical rigour --- survival curves, odds ratios,
Gini coefficients --- through a multi-layer stack of SQL, Python notebooks, and a
Streamlit executive dashboard deployed to Cloud Run.

% -----------------------------------------------------------------------------
\section{Business Context}
\label{sec:ecom-context}
% -----------------------------------------------------------------------------

\subsection{The 3\% Repeat Purchase Problem}

Olist is a Brazilian marketplace intermediary: sellers list products on Olist's
platform and Olist handles the customer-facing experience. The dataset covers
99,441 orders placed between September 2016 and October 2018. Of 93,358 unique
customers, only 2,801 --- roughly 3.0\% --- ever placed a second delivered order.

This figure is the analytical anchor of the entire project. It frames every
downstream analysis as a search for signal within extreme class imbalance.
Repeat customers are worth \textbf{1.92x} a one-time customer in \emph{total}
revenue (R\$309 against R\$161), but the reason is order count, not basket size:
they place roughly 2.1 orders. \emph{Per order} they spend \textbf{0.91x} as much
(R\$146 against R\$161), and their first order is about 10\% smaller than a
one-and-done customer's. The high spenders are disproportionately the ones who
never come back, which inverts the naive reading --- retention is valuable
because of frequency, not because returning customers are richer.

\begin{keyinsight}
A 3\% base rate means that standard accuracy metrics are meaningless for
classification tasks. Any model predicting ``never repeat'' would achieve 97\%
accuracy while being analytically useless. The project therefore avoids binary
classification in favour of survival analysis (which respects censoring) and
logistic regression odds ratios (which quantify relative risk between groups).
\end{keyinsight}

\subsection{The \texttt{customer\_id} vs.\ \texttt{customer\_unique\_id} Distinction}

Olist's schema contains two customer identifiers that are easily confused:

\begin{itemize}
  \item \texttt{customer\_id}: generated per order. A single person placing three
    orders will have three distinct \texttt{customer\_id} values.
  \item \texttt{customer\_unique\_id}: the true person-level identifier, stable
    across all orders.
\end{itemize}

Using \texttt{customer\_id} directly would cause three critical errors: (1) it
would inflate the apparent customer count, (2) every order would appear as a
first purchase, and (3) the repeat rate would compute to zero by construction.
All SQL scripts and Python code in this project join through
\texttt{customer\_id} only to resolve the link to \texttt{customer\_unique\_id},
then group exclusively on the latter.

\begin{lstlisting}[style=sql, caption={Wrong vs.\ correct customer grouping}]
-- WRONG: one row per order, not per person
SELECT customer_id, COUNT(*) AS orders
FROM olist_orders
GROUP BY customer_id;

-- CORRECT: one row per real customer
SELECT c.customer_unique_id, COUNT(*) AS orders
FROM olist_orders  o
JOIN olist_customers c ON o.customer_id = c.customer_id
GROUP BY c.customer_unique_id;
\end{lstlisting}

\subsection{Olist Marketplace Model}

Because Olist is an intermediary, the customer experience depends on three
parties: Olist (platform, review system, logistics coordination), the individual
seller (product quality, dispatch time), and the carrier. This creates a
confounding structure for causal claims: a poor delivery experience may reflect
the seller, the carrier, or geographic distance --- not Olist platform quality
per se. The analysis treats delivery days and review scores as observed outcomes
without attributing causality to any single party.

% -----------------------------------------------------------------------------
\section{Olist Dataset}
\label{sec:ecom-data}
% -----------------------------------------------------------------------------

\sourcefiles{%
  \filepath{streamlit/pages/5\_metodologia.py} (tab ``Fuente de Datos''),
  \filepath{notebooks/01\_data\_ingestion\_cleaning.ipynb}
}

\subsection{Nine CSV Files, One Analytical Table}

The raw dataset consists of nine CSV files joined into a single master table
(\filepath{data/processed/orders\_enriched.parquet}) during notebook 01. The
join graph is star-shaped: \texttt{olist\_orders} is the central fact table,
with dimension tables for customers, items, payments, reviews, products, sellers,
geolocation, and a category translation lookup.

\begin{table}[h]
\centering
\caption{Olist dataset files and approximate row counts}
\label{tab:olist-files}
\begin{tabular}{llr}
\toprule
\textbf{CSV file} & \textbf{Content} & \textbf{Rows (approx.)} \\
\midrule
\texttt{olist\_orders\_dataset.csv}         & Orders: status, timestamps, customer\_id & 99,441 \\
\texttt{olist\_order\_items\_dataset.csv}   & Line items: price, freight, seller       & 112,650 \\
\texttt{olist\_customers\_dataset.csv}      & customer\_id $\to$ unique\_id, state     & 99,441 \\
\texttt{olist\_products\_dataset.csv}       & Category, dimensions, weight             & 32,951 \\
\texttt{olist\_sellers\_dataset.csv}        & Seller state and city                    & 3,095 \\
\texttt{olist\_order\_payments\_dataset.csv}& Payment type, value, instalments         & 103,886 \\
\texttt{olist\_order\_reviews\_dataset.csv} & Review score 1--5, free text             & 100,000 \\
\texttt{olist\_geolocation\_dataset.csv}    & ZIP $\to$ lat/lon                        & 1,000,163 \\
\texttt{product\_category\_name\_translation.csv} & Portuguese $\to$ English category  & 71 \\
\bottomrule
\end{tabular}
\end{table}

The enriched parquet adds 30 derived columns including \texttt{delivery\_days},
\texttt{is\_late\_delivery}, \texttt{cohort\_month}, \texttt{months\_since\_cohort},
and aggregated payment totals. The analysis filters to 96,478 delivered orders
only; cancelled, unavailable, and invoiced orders are excluded because their
customer journey is incomplete.

\subsection{Right-Censoring Problem}

The observation window closes at October 2018. Any customer who made their first
purchase in, say, August 2018 has had at most two months to return. We cannot
observe whether they return in November 2018 or later. This is right-censoring
in the survival analysis sense: the event (second purchase) may have happened
after the observation period ended.

\begin{decisionbox}
Three approaches were considered for handling censoring:
\begin{enumerate}
  \item \textbf{Ignore it}: count non-repeaters as definitive churners. This
    underestimates the true repeat rate for recent cohorts.
  \item \textbf{Truncate}: only analyse cohorts old enough to have a full 12-month
    window. This wastes 40\%+ of the data and introduces selection bias toward early cohorts.
  \item \textbf{Kaplan-Meier} (chosen): treat non-repeaters as right-censored
    observations. This correctly accounts for the fact that recent customers have
    had less time to return, without discarding them.
\end{enumerate}
The Kaplan-Meier approach is implemented via the \texttt{lifelines} library.
The \texttt{event\_observed} column is 1 for confirmed repeaters and 0 for
right-censored customers.
\end{decisionbox}

% -----------------------------------------------------------------------------
\section{Cohort Retention}
\label{sec:ecom-cohorts}
% -----------------------------------------------------------------------------

\sourcefiles{%
  \filepath{sql/01\_cohort\_assignment.sql},
  \filepath{sql/02\_retention\_matrix.sql},
  \filepath{notebooks/03\_cohort\_retention.ipynb},
  \filepath{streamlit/pages/2\_retencion\_cohortes.py}
}

\subsection{Cohort Assignment Logic}

A cohort is defined by the month of a customer's first \emph{delivered} order.
The SQL implementation uses two CTEs: \texttt{customer\_first\_purchase} computes
the minimum \texttt{order\_purchase\_timestamp} per \texttt{customer\_unique\_id},
and \texttt{cohort\_assignment} computes \texttt{months\_since\_cohort} for every
subsequent order.

\begin{lstlisting}[style=sql, caption={Cohort assignment and tenure calculation (from \filepath{sql/01\_cohort\_assignment.sql})}]
customer_first_purchase AS (
    SELECT customer_unique_id,
           MIN(order_purchase_timestamp) AS first_purchase_date
    FROM delivered_orders
    GROUP BY customer_unique_id
),
cohort_assignment AS (
    SELECT d.customer_unique_id,
           DATE_TRUNC('month', fp.first_purchase_date) AS cohort_month,
           (EXTRACT(YEAR FROM AGE(
               DATE_TRUNC('month', d.order_purchase_timestamp),
               DATE_TRUNC('month', fp.first_purchase_date)
           )) * 12
           + EXTRACT(MONTH FROM AGE(
               DATE_TRUNC('month', d.order_purchase_timestamp),
               DATE_TRUNC('month', fp.first_purchase_date)
           )))::INT AS months_since_cohort
    FROM delivered_orders AS d
    JOIN customer_first_purchase AS fp
      ON d.customer_unique_id = fp.customer_unique_id
)
\end{lstlisting}

Note that \texttt{DATE\_TRUNC} is applied before \texttt{AGE} to ensure tenure
is measured in whole months, not fractional months. The final \texttt{SELECT}
also adds \texttt{ROW\_NUMBER() OVER (PARTITION BY customer\_unique\_id ORDER BY
order\_purchase\_timestamp)} to assign a sequential order number, which is used
downstream to isolate first orders in the activation analysis.

\subsection{Retention Matrix Construction}

The retention matrix pivots cohort-month pairs into a wide format using
\texttt{CASE-WHEN} expressions. Each cell reports the percentage of the cohort's
customers who placed a delivered order in that period:

\begin{equation}
\text{Retention}_{c,k} = \frac{\left|\{u \in C_c : \text{purchase in month } k\}\right|}{|C_c|}
\label{eq:retention}
\end{equation}

where $C_c$ is the set of customers in cohort $c$ and $k$ is the number of months
since cohort formation. By construction, $\text{Retention}_{c,0} = 100\%$ for
all cohorts.

\begin{lstlisting}[style=sql, caption={Retention matrix pivot (excerpt from \filepath{sql/02\_retention\_matrix.sql})}]
SELECT
    s.cohort_month,
    s.cohort_size,
    ROUND(MAX(CASE WHEN m.months_since_cohort = 0
        THEN m.active_customers * 100.0 / s.cohort_size END), 2) AS month_0,
    ROUND(MAX(CASE WHEN m.months_since_cohort = 1
        THEN m.active_customers * 100.0 / s.cohort_size END), 2) AS month_1,
    ROUND(MAX(CASE WHEN m.months_since_cohort = 2
        THEN m.active_customers * 100.0 / s.cohort_size END), 2) AS month_2
    -- ... through month_11
FROM cohort_sizes AS s
LEFT JOIN monthly_activity AS m ON s.cohort_month = m.cohort_month
GROUP BY s.cohort_month, s.cohort_size
ORDER BY s.cohort_month;
\end{lstlisting}

\subsection{Count-Based vs.\ Revenue-Based Retention}

The Streamlit dashboard (page 2) offers a toggle between two retention modes:

\begin{itemize}
  \item \textbf{Count-based}: the number of distinct customers active in month $k$
    divided by cohort size. Stored in \filepath{data/processed/cohort\_retention\_matrix.parquet}.
  \item \textbf{Revenue-based}: the total payment value from returning customers in
    month $k$ divided by month-0 revenue. Stored in
    \filepath{data/processed/cohort\_revenue\_retention.parquet}.
\end{itemize}

Revenue retention amplifies the signal of high-value returners. A cohort where
one high-spend customer returns could show high revenue retention but low count
retention. The divergence between the two modes is itself informative: large
gaps suggest that repeat purchases are concentrated among a few high-value
customers rather than broadly distributed.

The Streamlit code normalises both matrices identically:
\begin{lstlisting}[style=python, caption={Normalisation to retention percentages (\filepath{streamlit/pages/2\_retencion\_cohortes.py}, line~90)}]
data_pct = data.div(data.iloc[:, 0], axis=0) * 100
\end{lstlisting}

% -----------------------------------------------------------------------------
\section{Kaplan-Meier Survival Analysis}
\label{sec:ecom-km}
% -----------------------------------------------------------------------------

\sourcefiles{%
  \filepath{notebooks/03\_cohort\_retention.ipynb},
  \filepath{streamlit/pages/2\_retencion\_cohortes.py} (tab ``Kaplan-Meier''),
  \filepath{streamlit/pages/5\_metodologia.py} (tab ``Supervivencia Kaplan-Meier'')
}

\subsection{Estimator Definition}

The Kaplan-Meier (KM) estimator computes the probability that the event of
interest --- here, a second purchase --- has \emph{not} yet occurred by time $t$.
In this project, $S(t)$ is the probability that a customer has not returned
within $t$ days of their first purchase:

\begin{equation}
\hat{S}(t) = \prod_{t_i \leq t} \left(1 - \frac{d_i}{n_i}\right)
\label{eq:km}
\end{equation}

where:
\begin{itemize}
  \item $t_i$ are the ordered distinct event times (days when second purchases occurred),
  \item $d_i$ is the number of second purchases at time $t_i$ (``deaths''),
  \item $n_i$ is the number of customers still at risk at time $t_i$ (neither
    purchased again nor censored before $t_i$).
\end{itemize}

\subsection{Preparing the Survival Dataset}

The \texttt{survival\_data.parquet} file contains one row per customer with two
key columns:

\begin{lstlisting}[style=python, caption={Survival data construction (from \filepath{streamlit/pages/5\_metodologia.py})}]
survival["duration_days"] = np.where(
    survival["is_repeat_customer"],
    (survival["second_purchase_date"] - survival["first_purchase_date"]).dt.days,
    (DATASET_END_DATE - survival["first_purchase_date"]).dt.days,
)
survival["event_observed"] = survival["is_repeat_customer"].astype(int)
\end{lstlisting}

For the 97\% of customers who never returned, \texttt{event\_observed = 0}
and \texttt{duration\_days} measures how long we watched them without seeing a
repeat purchase --- the censoring time, not the time-to-event.

\subsection{Fitting the Model in Streamlit}

The dashboard fits a \texttt{KaplanMeierFitter} from the \texttt{lifelines}
library (version \texttt{>=0.29}) interactively:

\begin{lstlisting}[style=python, caption={KM fitting, originally in \filepath{streamlit/pages/2\_retencion\_cohortes.py}; now precomputed by \filepath{data-pipeline/05\_build\_web\_json.py}}]
from lifelines import KaplanMeierFitter

kmf = KaplanMeierFitter()
kmf.fit(
    durations=survival_clean["duration_days"],
    event_observed=survival_clean["event_observed"],
)
km_df = kmf.survival_function_.reset_index()
ci = kmf.confidence_interval_survival_function_.reset_index()
median_surv = kmf.median_survival_time_   # returns inf: S(t) never reaches 0.5
\end{lstlisting}

\begin{keyinsight}
That last line is a trap worth naming. \texttt{median\_survival\_time\_} returns
\texttt{inf} on this dataset, because $\hat{S}(t)$ plateaus near 0.95 and never
crosses 0.5. Printing it as ``the median'' produces a headline figure that does
not exist. The quotable statistic is conditional on returning, and has to be
computed separately over the uncensored rows.
\end{keyinsight}

The dashboard also supports segmented KM curves: the user selects
\texttt{payment\_type} or \texttt{customer\_state} as a stratification variable,
and the page fits a separate KM curve per group. This directly supports the
log-rank test finding that voucher payers have statistically different survival
curves from credit card and boleto payers.

\subsection{Interpreting the Curve in Context}

For Olist, $\hat{S}(t)$ starts at 1.0 (all customers are ``surviving'' ---
i.e., have not yet reordered) and decays slowly, plateauing near 0.95. The decay
\emph{decelerates but never stops}: the curve loses 1.23 percentage points in the
first quarter and roughly 0.6 points in every quarter thereafter, out to two
years.

\begin{keyinsight}
This corrects an earlier reading of this same dataset, which described reorder
probability as ``practically nil'' after six months. The curve does not support
that. Repurchase keeps trickling in for the full two-year observation window; it
merely slows. The distinction matters for the recommendation, because a
six-month cutoff would write off customers who are still converting.
\end{keyinsight}

Because $\hat{S}(t)$ never crosses 0.5, \textbf{the median survival time does not
exist} --- it is undefined, not merely uninterpretable. What can be quoted is a
conditional statistic: among the 3\% who \emph{did} return, the median wait was
\textbf{82 days}. That is a different quantity and must be labelled as such.

\begin{keyinsight}
The KM curve shows where the density is, not where the opportunity ends. Among
customers who returned, the median wait was 82 days and 27\% came back within
30 --- so the first quarter is the right place to concentrate an intervention.
But since the curve keeps declining for two years, a later touch is not wasted
effort, and framing day 90 as a cutoff would be a misreading of the same data
that motivates the campaign.
\end{keyinsight}

% -----------------------------------------------------------------------------
\section{RFM Segmentation}
\label{sec:ecom-rfm}
% -----------------------------------------------------------------------------

\sourcefiles{%
  \filepath{sql/03\_rfm\_segmentation.sql},
  \filepath{notebooks/04\_rfm\_ltv\_activation.ipynb},
  \filepath{streamlit/pages/3\_segmentos\_clientes.py} (tabs ``Distribucion'' through ``Lorenz'')
}

\subsection{RFM Dimensions and Quintile Scoring}

RFM scores every customer on three behavioural dimensions relative to the
analysis reference date (one day after the last observed purchase):

\begin{table}[h]
\centering
\caption{RFM dimension definitions and scoring direction}
\label{tab:rfm-dims}
\begin{tabular}{llll}
\toprule
\textbf{Dimension} & \textbf{Metric} & \textbf{Direction} & \textbf{Score 5 means} \\
\midrule
Recency (R)   & Days since last purchase  & Lower is better  & Most recent 20\%  \\
Frequency (F) & Count of distinct orders  & Higher is better & Top 20\% by orders \\
Monetary (M)  & Total lifetime spend (R\$) & Higher is better & Top 20\% by revenue \\
\bottomrule
\end{tabular}
\end{table}

Quintile scoring uses \texttt{NTILE(5)} in PostgreSQL. Recency requires reversed
ordering because fewer days since last purchase is better:

\begin{lstlisting}[style=sql, caption={Quintile scoring in \filepath{sql/03\_rfm\_segmentation.sql}}]
NTILE(5) OVER (ORDER BY recency_days DESC)  AS r_score,  -- 5=most recent
NTILE(5) OVER (ORDER BY frequency ASC)      AS f_score,
NTILE(5) OVER (ORDER BY monetary ASC)       AS m_score
\end{lstlisting}

\subsection{Seven Segment Definitions}

Score combinations are mapped to seven business segments through a \texttt{CASE}
statement. The segment rules below are implemented identically in both SQL
(\filepath{sql/03\_rfm\_segmentation.sql}) and Python
(\filepath{streamlit/pages/5\_metodologia.py}):

\begin{table}[h]
\centering
\caption{RFM segment assignment rules}
\label{tab:rfm-segments}
\begin{tabular}{llll}
\toprule
\textbf{Segment (English)} & \textbf{R score} & \textbf{F score} & \textbf{M score} \\
\midrule
Champions                 & $\geq 4$  & $\geq 4$  & $\geq 4$  \\
Loyal Customers           & any       & $\geq 4$  & $\geq 3$  \\
Potential Loyalists       & $\geq 4$  & 2--3      & any       \\
New Customers             & $\geq 4$  & 1         & any       \\
At Risk                   & 2--3      & $\geq 3$  & any       \\
Hibernating               & 2--3      & 1--2      & any       \\
Lost                      & 1         & any       & any       \\
\bottomrule
\end{tabular}
\end{table}

\subsection{Why RFM over K-Means for This Use Case}

\begin{decisionbox}
K-Means clustering was considered as an alternative segmentation approach.
Three arguments favour RFM quintiles for this dataset:
\begin{enumerate}
  \item \textbf{Interpretability}: Business stakeholders can understand and act
    on ``Champions'' and ``At Risk'' segments. K-Means cluster labels are
    arbitrary without post-hoc labelling.
  \item \textbf{Class imbalance}: With 97\% single-purchase customers, K-Means
    on raw RFM values would produce clusters dominated by F=1 customers with
    little sub-group differentiation. Quintile scoring spreads customers across
    scores regardless of the raw distribution.
  \item \textbf{Operational directness}: Each segment maps to a specific
    marketing action (reactivation for ``At Risk'', retention for ``Champions'').
    K-Means clusters require interpretation before action; RFM segments are
    pre-labelled for action.
\end{enumerate}
The downside of NTILE-based RFM is that quintile boundaries are not stable: they
shift if the customer base grows. For production use, fixed absolute thresholds
(e.g., R$<$30 days = score 5) would be more stable but require domain calibration.
\end{decisionbox}

% -----------------------------------------------------------------------------
\section{Activation Analysis}
\label{sec:ecom-activation}
% -----------------------------------------------------------------------------

\sourcefiles{%
  \filepath{sql/04\_ltv\_activation.sql} (Part 2),
  \filepath{notebooks/04\_rfm\_ltv\_activation.ipynb},
  \filepath{streamlit/pages/3\_segmentos\_clientes.py} (tab ``Activacion'')
}

\subsection{Logistic Regression Setup}

The activation analysis estimates the association between first-order
characteristics and the probability of making a repeat purchase. The binary
outcome is \texttt{is\_repeat} (1 if two or more delivered orders, 0 otherwise).
The model is a standard logistic regression:

\begin{equation}
\log\!\left(\frac{p}{1-p}\right) = \beta_0 + \beta_1 x_1 + \cdots + \beta_k x_k
\label{eq:logistic}
\end{equation}

where $p = \Prob(\text{repeat} = 1 \mid \mathbf{x})$ and predictors $x_i$
include payment type (dummy-coded), first-order value (continuous), item count
(categorical), and product category (top categories dummy-coded). The model is
fitted with \texttt{statsmodels} (\texttt{Logit} class) to obtain coefficient
estimates with standard errors and Wald test $p$-values.

\subsection{Odds Ratios and Their Interpretation}

Each logistic regression coefficient $\hat{\beta}_i$ is exponentiated to yield
an odds ratio:

\begin{equation}
\widehat{\text{OR}}_i = \exp(\hat{\beta}_i)
\label{eq:odds-ratio}
\end{equation}

An odds ratio greater than 1 indicates that the feature is associated with
higher odds of repeat purchase, holding other features constant. The stored
artefact \filepath{data/processed/activation\_coefficients.parquet} contains
columns \texttt{feature}, \texttt{odds\_ratio}, \texttt{p\_value},
\texttt{ci\_lower}, and \texttt{ci\_upper} (95\% Wald confidence intervals on
the log-odds scale, exponentiated).

\begin{keyinsight}
\textbf{Voucher OR = 1.42}: customers who paid their first order with a voucher
(gift card or discount code) have 1.42 times the odds of returning compared to
other payment methods, after controlling for order value, item count, and
category. This is the strongest statistically significant predictor in the model.

Business interpretation: voucher users may have lower price sensitivity or may
have been acquired through promotional channels that attract higher-intent buyers.
The recommendation is to expand voucher/discount programmes targeting first-time
buyers.
\end{keyinsight}

The dashboard renders the odds ratios on a log$_2$ scale (forest plot style)
so that symmetric bar lengths represent equivalent fold-changes above and below
the null (OR = 1). Green bars indicate significantly positive associations
($p < 0.05$, OR $> 1$); red bars indicate significantly negative associations.

\subsection{SQL Activation Rate Table}

The SQL script \filepath{sql/04\_ltv\_activation.sql} (Part 2) computes raw
repeat rates by feature group --- a fast univariate view complementing the
multivariate logistic regression. It uses a \texttt{UNION ALL} to stack results
for payment type, item count bucket, and order value quartile:

\begin{lstlisting}[style=sql, caption={Activation rates by payment type (excerpt from \filepath{sql/04\_ltv\_activation.sql})}]
SELECT
    'payment_type'                       AS feature,
    first_payment_type                   AS feature_value,
    COUNT(*)                             AS customers,
    SUM(is_repeat)                       AS repeat_customers,
    ROUND(100.0 * SUM(is_repeat) / COUNT(*), 2) AS repeat_rate_pct
FROM first_order_features
GROUP BY first_payment_type
\end{lstlisting}

% -----------------------------------------------------------------------------
\section{Revenue Concentration: Lorenz Curve and Gini Coefficient}
\label{sec:ecom-lorenz}
% -----------------------------------------------------------------------------

\sourcefiles{%
  \filepath{streamlit/pages/3\_segmentos\_clientes.py} (tab ``Lorenz''),
  \filepath{notebooks/04\_rfm\_ltv\_activation.ipynb}
}

\subsection{Lorenz Curve Construction}

The Lorenz curve plots cumulative revenue share against cumulative customer
share, both sorted from lowest to highest spender. At point $(p, L(p))$, the
bottom $p$ fraction of customers account for $L(p)$ fraction of total revenue.
Perfect equality would yield $L(p) = p$ (the diagonal).

The Python implementation in \filepath{streamlit/pages/3\_segmentos\_clientes.py}
computes the curve directly from the \texttt{total\_revenue} column of
\texttt{rfm\_segments.parquet}:

\begin{lstlisting}[style=python, caption={Lorenz curve and Gini calculation (\filepath{streamlit/pages/3\_segmentos\_clientes.py}, lines~329--335)}]
revenue_sorted = np.sort(rfm_raw["total_revenue"].values)
n = len(revenue_sorted)
cum_revenue = np.cumsum(revenue_sorted) / revenue_sorted.sum()
cum_population = np.arange(1, n + 1) / n

gini = 1 - 2 * np.trapezoid(cum_revenue, cum_population)
\end{lstlisting}

\subsection{Gini Coefficient Formula}

The Gini coefficient is twice the area between the Lorenz curve and the
line of equality:

\begin{equation}
G = 1 - 2 \int_0^1 L(p)\,\mathrm{d}p
\label{eq:gini}
\end{equation}

In the discrete implementation, the integral is approximated with the trapezoidal
rule (\texttt{np.trapezoid} in NumPy $\geq$ 2.0, falling back to \texttt{np.trapz}
for earlier versions). $G = 0$ means perfect equality; $G = 1$ means one customer
holds all revenue.

\subsection{Business Interpretation}

The Olist customer revenue distribution is highly concentrated. Quantifying the
top-20\% contribution directly:

\begin{lstlisting}[style=python, caption={Top-20\% revenue share calculation}]
top_20_rev = cum_revenue[int(0.8 * n)]
top_20_share = (1 - top_20_rev) * 100  # % held by top 20%
\end{lstlisting}

The Champions segment (R $\geq 4$, F $\geq 4$, M $\geq 4$) --- approximately
0.1\% of customers --- generates outsized lifetime value (R\$\,545 per customer
at 12 months vs.\ R\$\,160 average). This concentration justifies a differentiated
retention programme for the top revenue tier.

\subsection{LTV Curves by Segment}

Cumulative LTV is computed in \filepath{sql/04\_ltv\_activation.sql} (Part 1)
using a running sum window function:

\begin{lstlisting}[style=sql, caption={Cumulative LTV per customer (from \filepath{sql/04\_ltv\_activation.sql})}]
SUM(cr.monthly_revenue) OVER (
    PARTITION BY cr.cohort_month
    ORDER BY cr.months_since_cohort
    ROWS UNBOUNDED PRECEDING
) / cs.cohort_size  AS ltv_per_customer
\end{lstlisting}

The Streamlit page loads \filepath{data/processed/ltv\_curves.parquet} via
\texttt{load\_ltv\_curves()} and renders a multi-line chart with one series per
RFM segment, coloured by the shared categorical palette from
\filepath{streamlit/utils/charts.py}.

% -----------------------------------------------------------------------------
\section{Dashboard Architecture: From Streamlit to a Static Export}
\label{sec:ecom-streamlit}
% -----------------------------------------------------------------------------

\sourcefiles{%
  \filepath{streamlit/app.py},
  \filepath{streamlit/utils/data\_loader.py},
  \filepath{streamlit/utils/filters.py},
  \filepath{streamlit/pages/} (6 page files)
}

This dashboard was built twice, and the reason for the rebuild is the most
transferable thing in the chapter. The first implementation was a six-page
Streamlit app on Cloud Run; the second is a six-page Next.js static export served
from the portfolio hub at \texttt{data-analyst.gonor.me/cohorts}, with no backend
at all. Both sections below are kept: the Streamlit design is described first
because the rebuild only became visible \emph{through} it.

\subsection{First Implementation: Streamlit Page Structure}

The original dashboard was a six-page Streamlit multi-page app launched via:
\begin{lstlisting}[style=dockerfile, caption={Run command (from \filepath{Dockerfile})}]
CMD ["streamlit", "run", "streamlit/app.py",
     "--server.port=8501", "--server.address=0.0.0.0",
     "--server.headless=true"]
\end{lstlisting}

\begin{table}[h]
\centering
\caption{Streamlit pages and their analytical content}
\label{tab:streamlit-pages}
\begin{tabular}{lll}
\toprule
\textbf{Page file} & \textbf{Title} & \textbf{Key content} \\
\midrule
\texttt{1\_resumen\_ejecutivo.py}   & Resumen Ejecutivo    & KPI cards, revenue trend, retention funnel \\
\texttt{2\_retencion\_cohortes.py}  & Retención por Cohortes & Heatmaps, KM survival, cohort comparator \\
\texttt{3\_segmentos\_clientes.py}  & Segmentos de Clientes & RFM scatter, LTV, activation OR, Lorenz \\
\texttt{4\_analisis\_geografico.py} & Análisis Geográfico  & State rankings, delivery-retention scatter \\
\texttt{5\_metodologia.py}          & Metodología          & Definitions, calculation logic, data docs \\
\texttt{6\_proceso\_tecnico.py}     & Proceso Técnico      & Embedded Jupyter notebooks as HTML iframes \\
\bottomrule
\end{tabular}
\end{table}

\subsection{Data Loading and Caching (Streamlit)}

All data loading was centralised in \filepath{streamlit/utils/data\_loader.py}.
Each parquet file has a dedicated \texttt{@st.cache\_data} function so that
repeated page navigations do not trigger re-reads from disk. The eight loader
functions are:

\begin{lstlisting}[style=python, caption={Cached loaders in \filepath{streamlit/utils/data\_loader.py}}]
@st.cache_data
def load_orders()          -> pd.DataFrame | None: ...  # orders_enriched.parquet
def load_customers()       -> pd.DataFrame | None: ...  # customers_summary.parquet
def load_cohort_retention()-> pd.DataFrame | None: ...  # cohort_retention_matrix.parquet
def load_cohort_revenue()  -> pd.DataFrame | None: ...  # cohort_revenue_retention.parquet
def load_survival()        -> pd.DataFrame | None: ...  # survival_data.parquet
def load_rfm()             -> pd.DataFrame | None: ...  # rfm_segments.parquet
def load_ltv_curves()      -> pd.DataFrame | None: ...  # ltv_curves.parquet
def load_activation()      -> pd.DataFrame | None: ...  # activation_coefficients.parquet
\end{lstlisting}

Functions \texttt{load\_rfm()} and \texttt{load\_ltv\_curves()} additionally
apply \texttt{\_translate\_segments()} to map English segment names to their
Spanish equivalents for display (e.g., ``Champions'' $\to$ ``Alto Valor'').

\subsection{Global Filter System --- and What It Revealed}

Filters were maintained in \texttt{st.session\_state} and applied by helper
functions in \filepath{streamlit/utils/filters.py}. Three global filters are
available:

\begin{enumerate}
  \item \textbf{Date range} (\texttt{date\_start} / \texttt{date\_end}): applied
    to orders via \texttt{apply\_date\_filter()} and to customers via
    \texttt{apply\_date\_filter\_customers()} (which filters on
    \texttt{cohort\_month} rather than \texttt{order\_purchase\_timestamp}).
  \item \textbf{Minimum cohort size} (\texttt{min\_cohort\_size}, default 50):
    \texttt{apply\_cohort\_size\_filter()} removes cohorts with fewer than
    \texttt{min\_cohort\_size} customers in month 0 from the retention matrices.
  \item \textbf{RFM segment selection} (\texttt{selected\_segments}):
    \texttt{apply\_segment\_filter()} subsets the RFM dataframe to selected segments.
\end{enumerate}

Active filters are shown via \texttt{render\_active\_filter\_badges()}, which
renders coloured HTML badges without recomputing any data.

\begin{keyinsight}
Read those three filters again and note what none of them does: recompute
anything from the 96,478 delivered orders. A date range \emph{selects columns} of
an already-aggregated retention matrix; a minimum cohort size \emph{drops rows}
from it; a segment selection \emph{subsets} an RFM table that was scored once,
offline. Every one is a projection over a precomputed aggregate.

That is the entire argument for the rebuild. A stateful Python server is only
required when a user action changes the computation. Here it never did --- so the
server was doing nothing that a browser could not do over a small JSON payload,
while costing a container, a cold start, and a second design language on the
same domain.
\end{keyinsight}

\subsection{Second Implementation: Zero-Backend Static Export}

The rebuild pre-aggregates 36\,MB of parquet into \textbf{276\,KB of JSON}
(roughly 38\,KB gzipped) via \filepath{data-pipeline/05\_build\_web\_json.py}, and
ships a Next.js application compiled with \texttt{output: 'export'}. The six
pages, the filters and the interactivity all survive; the Cloud Run service does
not.

\begin{table}[h]
\centering
\caption{What moved where in the rebuild}
\label{tab:ecom-rebuild}
\begin{tabular}{lll}
\toprule
\textbf{Concern} & \textbf{Streamlit} & \textbf{Static export} \\
\midrule
Aggregation      & per request, in Python & once, in the build pipeline \\
Filter state     & \texttt{st.session\_state} & React state in the browser \\
Kaplan--Meier    & fitted per page load & precomputed, shipped as points \\
Transport        & Cloud Run HTTP & static JSON under \texttt{public/cohorts/data/} \\
Notebooks        & \texttt{components.v1.html} iframe & same HTML, served statically \\
Cold start       & $\sim$10\,s & none \\
\bottomrule
\end{tabular}
\end{table}

Two consequences are worth stating because they generalise beyond this project:

\begin{enumerate}
  \item \textbf{Correctness has to be proved, not assumed.} Moving the
    computation from request time to build time means a silent divergence between
    the parquet and the published JSON would be invisible.
    \filepath{data-pipeline/06\_verify\_parity.py} re-derives every published
    figure from the parquet and checks the Kaplan--Meier curve and both
    confidence bounds against \texttt{lifelines}. It cannot run in CI, because
    the parquet lives only in GCS.
  \item \textbf{The data became a build input, not an artifact.} The repository's
    global \texttt{*.json} ignore rule excluded the very files the dashboard
    renders from. The build still succeeded and would have deployed a dashboard
    with no numbers on it --- a failure mode that returns HTTP 200. It is now
    covered by an explicit negation in \filepath{.gitignore}, a gate test
    asserting a rendered figure, and a health probe on
    \texttt{/cohorts/data/meta.json}.
\end{enumerate}

\subsection{Publishing the Notebooks (``Proceso Técnico'')}

Page 6 publishes the analytical pipeline inside the dashboard itself. The four
Jupyter notebooks are pre-converted to static HTML using:

\begin{lstlisting}[style=dockerfile, caption={Notebook export command}]
jupyter nbconvert --to html \
  --output-dir=streamlit/notebooks_html \
  notebooks/*.ipynb
\end{lstlisting}

The HTML files are committed to git. Under Streamlit they were
\texttt{COPY}-ed into the Docker image and rendered with
\texttt{streamlit.components.v1.html()}:

\begin{lstlisting}[style=python, caption={Notebook rendering in \filepath{streamlit/pages/6\_proceso\_tecnico.py}, lines~113--116}]
with open(path, "r", encoding="utf-8") as f:
    html_content = f.read()

components.html(html_content, height=800, scrolling=True)
\end{lstlisting}

The static export keeps the pattern and drops the container: the same HTML now
lives under the app's own \texttt{public/} slug and is rendered in a plain
\texttt{<iframe>}. Either way, each tab shows one notebook with a description
card explaining its inputs and outputs, giving portfolio reviewers full
visibility into the analytical pipeline without leaving the dashboard.

\begin{keyinsight}
Assets must live under the application's own slug
(\texttt{public/cohorts/\ldots}), never at the root. Seven dashboards share one
origin in the merged hub, and a root-level asset from one of them would collide
with another's. See Chapter~\ref{ch:infrastructure}.
\end{keyinsight}

% -----------------------------------------------------------------------------
\section{SQL Scripts Reference}
\label{sec:ecom-sql}
% -----------------------------------------------------------------------------

\sourcefiles{%
  \filepath{sql/01\_cohort\_assignment.sql},
  \filepath{sql/02\_retention\_matrix.sql},
  \filepath{sql/03\_rfm\_segmentation.sql},
  \filepath{sql/04\_ltv\_activation.sql},
  \filepath{sql/05\_geographic\_retention.sql}
}

The five SQL scripts are written for PostgreSQL and share a common structure:
delivered orders are always the base CTE (filtering \texttt{order\_status =
'delivered'}), and \texttt{customer\_unique\_id} is consistently used as the
person-level key.

\begin{table}[h]
\centering
\caption{SQL scripts, key techniques, and business questions}
\label{tab:sql-scripts}
\begin{tabularx}{\textwidth}{lXl}
\toprule
\textbf{Script} & \textbf{Business question answered} & \textbf{Key technique} \\
\midrule
\texttt{01\_cohort\_assignment.sql} &
  Which cohort does each customer belong to? &
  \texttt{ROW\_NUMBER()}, \texttt{DATE\_TRUNC}, \texttt{AGE} \\
\texttt{02\_retention\_matrix.sql} &
  What \% of each cohort returns each month? &
  \texttt{CASE-WHEN} pivot, \texttt{COUNT(DISTINCT)} \\
\texttt{03\_rfm\_segmentation.sql} &
  How do customers segment by RFM? &
  \texttt{NTILE(5)}, \texttt{CASE} segment mapping \\
\texttt{04\_ltv\_activation.sql} &
  What drives repeat purchase (activation)? &
  Running \texttt{SUM() OVER}, \texttt{DISTINCT ON}, \texttt{UNION ALL} \\
\texttt{05\_geographic\_retention.sql} &
  Which states retain better and why? &
  \texttt{FILTER (WHERE ...)}, \texttt{RANK() OVER} \\
\bottomrule
\end{tabularx}
\end{table}

The geographic retention script uses the \texttt{FILTER} aggregate modifier
to compute on-time delivery rate and repeat customer count in a single pass:

\begin{lstlisting}[style=sql, caption={Conditional aggregation with FILTER (from \filepath{sql/05\_geographic\_retention.sql})}]
ROUND(100.0 * COUNT(*) FILTER (
    WHERE d.order_delivered_customer_date <= d.order_estimated_delivery_date
) / COUNT(*), 2)                           AS on_time_pct,

COUNT(DISTINCT customer_unique_id) FILTER (
    WHERE order_count >= 2
)                                          AS repeat_customers
\end{lstlisting}

States with fewer than 50 customers are filtered out of the ranking to avoid
noise from low sample sizes (\texttt{WHERE total\_customers >= 50}).

% -----------------------------------------------------------------------------
\section{Key Findings Summary}
\label{sec:ecom-findings}
% -----------------------------------------------------------------------------

\begin{table}[h]
\centering
\caption{Key quantitative findings from the Olist analysis}
\label{tab:ecom-findings}
\begin{tabular}{lll}
\toprule
\textbf{Finding} & \textbf{Value} & \textbf{Implication} \\
\midrule
Repeat purchase rate          & 3.0\%   & Extreme retention challenge \\
Repeat customer total revenue & 1.92x   & Driven by 2.1 orders, not basket size \\
Repeat customer \emph{per order}  & 0.91x   & Returning buyers spend \emph{less} per order \\
Voucher payment OR            & 1.42    & Expand voucher first-order incentives \\
Late delivery rate            & 8.1\%   & Delivery accuracy is a retention lever \\
Median wait among returners   & 82 days & Concentrate, do not cut off, at Q1 \\
Champions LTV (12 months)     & R\$\,545 vs.\ R\$\,160 avg & Prioritise top segment \\
\bottomrule
\end{tabular}
\end{table}

% -----------------------------------------------------------------------------
\section{Challenge and Interview Questions}
\label{sec:ecom-questions}
% -----------------------------------------------------------------------------

\begin{challengebox}
\textbf{Retention matrix normalisation.} The retention matrix is constructed by
dividing each row by its month-0 count. What happens if a customer places two
orders in month 0? Are they counted twice in the denominator? Trace through the
SQL in \texttt{02\_retention\_matrix.sql} to determine how this is handled.
\end{challengebox}

\begin{challengebox}
\textbf{Censoring bias direction.} The dataset ends in October 2018. For cohorts
formed in September 2018, customers have had at most two months to reorder. In
which direction does ignoring censoring bias the estimated repeat rate --- and
for which cohorts is the bias largest?
\end{challengebox}

\begin{challengebox}
\textbf{Gini coefficient edge case.} If all customers have exactly equal revenue,
$G = 0$. If one customer holds all revenue, $G = 1$. Verify that the discrete
trapezoidal approximation in the code converges to these limits for $n = 100$
customers.
\end{challengebox}

\begin{challengebox}
\textbf{NTILE stability.} The RFM scores are computed as quintiles of the current
customer base. If 10,000 new customers join with low recency scores, how do the
quintile boundaries shift, and what happens to existing Champions?
\end{challengebox}

\begin{interviewbox}
\textbf{Why Kaplan-Meier rather than a simple proportion?} A hiring manager asks:
``Why not just compute repeat rate as 2801 / 93358 = 3\%?'' Explain in two sentences
what information the KM curve adds that the simple proportion does not.
\end{interviewbox}

\begin{interviewbox}
\textbf{Interpreting OR = 1.42.} Explain to a non-technical product manager what
it means that voucher payment has an odds ratio of 1.42 for repeat purchase.
Avoid the word ``odds''. Frame it in terms of risk or probability.
\end{interviewbox}

\begin{interviewbox}
\textbf{Confounding in the activation analysis.} A colleague points out that
voucher users may also have placed higher-value first orders, and higher order
value is itself associated with return probability. How does the multivariate
logistic regression address this concern, and what residual confounding might
remain?
\end{interviewbox}

\begin{interviewbox}
\textbf{Lorenz curve business case.} The Gini coefficient shows high revenue
concentration. A VP asks: ``Should we focus on acquiring more Champions or on
converting Hibernating customers?'' Walk through the data-driven argument for
each position and state which you would recommend.
\end{interviewbox}
